% =====================================================================
% Section: Experiments -- Qualitative (visual) comparison (outline.md §3.4;
%   defines sec:visual). Discusses Fig.7 (fig:visual).
% Draft v1 (2026-06-27). Samples chosen by GT high-frequency-energy ranking +
%   SR-literature convention (paper/scripts/build_fig7_assets.py). PSNR/SSIM are
%   read VERBATIM from paper/figures/fig7_assets/metrics.csv (Y channel, crop_border=4,
%   recomputed from published 8-bit PNGs; ~0.07 dB below the logged float means due to
%   PNG quantization -- relative Bicubic-vs-DPRNet gaps are exact).
% HONEST BOUNDARY: competitor SR crops (original CATANet / SRFormer-light) are not yet
%   available locally; Fig.7's competitor column is a placeholder and this text makes
%   no DPRNet-vs-competitor visual claim until those crops are produced on the
%   training machine. Only DPRNet-vs-Bicubic / vs-GT statements are made here.
% =====================================================================
\subsection{Qualitative comparison}
\label{sec:visual}

Beyond aggregate PSNR/SSIM, Fig.~\ref{fig:visual} compares reconstructions at
$\times4$ on three texture-dense hard samples selected by ground-truth
high-frequency energy: two Urban100 images dominated by repetitive building
grids (\texttt{img\_092}, \texttt{img\_024}) and one Manga109 page of
high-contrast line art (\texttt{ThatsIzumiko}). These are exactly the regimes in
which content-based routing is expected to help most, since distant tokens that
share the same structure can be aggregated together rather than only within a
local window.

On the repetitive facades of Urban100, the bicubic baseline collapses adjacent
window edges into blurred, aliased bands, whereas DPRNet recovers the grid
spacing with markedly sharper boundaries; the per-crop scores rise from
$16.58$/$0.4377$ to $19.62$/$0.6922$ on \texttt{img\_092} and from
$18.14$/$0.4593$ to $21.05$/$0.7071$ on \texttt{img\_024}. On the Manga109 line
art, where thin strokes are easily broken by upsampling, DPRNet keeps the
contours continuous and improves the score from $18.00$/$0.5118$ to
$21.42$/$0.7597$. These observations are consistent with the routing diagnostics
of Sec.~\ref{sec:routing_analysis}: confidence-aware aggregation concentrates on
the structured, high-frequency regions that bicubic interpolation cannot resolve.
A side-by-side comparison against the strongest competitors
(CATANet~\cite{liu2025catanet}, SRFormer-light~\cite{zhou2023srformer}) on the
same crops is included in the final version.
